../cictikz/src/cictikz/data/tex/cictikz_preamble.tex

% The same one-bit string DAC as dac_r_div, but with two resistors instead
% of three, so the taps are V_REF itself and V_REF/2 rather than 2/3 and 1/3.
% The pair is the whole point of the slide: which of the two is the right
% way to build a DAC is the question the lecture asks next.

\begin{circuitikz}[american, thick, transform shape, circuit ee IEC]

  ../cictikz/src/cictikz/data/tex/cictikz_lib.tex
  ../cictikz/src/cictikz/data/tex/rdac_lib.tex

  \def\xs{0}
  \def\xo{3.2}

  % ---- Two resistors: the top tap is V_REF, the lower one is V_REF/2 ----
  \draw (\xs,0.5) \vresistor{$R$} -- ++(0,0.5) coordinate (tb)
                  \vresistor{$R$} -- ++(0,0.5) coordinate (ta);
  \draw (\xs,0.5) -- (\xs,0.2) \vground;

  \draw (ta) -- ++(-1.0,0) coordinate (vrefport);
  \rdacPort{(vrefport)}{south}{$V_{REF}$}

  % ---- b_0 passes V_REF, b_0-bar the midpoint ----
  \rdacSw{MA}{(ta)}{2.0}{$b_0$}
  \rdacSw{MB}{(tb)}{2.0}{$\overline{b_0}$}

  \coordinate (ra) at ($(ta)+(2.0,0)$);
  \coordinate (rb) at ($(tb)+(2.0,0)$);
  \draw (ra) -- (\xo,0 |- ra) -- (\xo,0 |- rb) -- (rb);
  \coordinate (vo) at ($(\xo,0 |- ra)!0.5!(\xo,0 |- rb)$);
  \fill (vo) circle (0.07);
  \draw (vo) -- ++(0.8,0) coordinate (voutport);
  \rdacPort{(voutport)}{south west}{$V_{OUT}$}

\end{circuitikz}
\end{document}
